% GENERATED FILE. Edit canonical lessons, then run npm run generate:books.
% canonical-source-hash: fnv1a64:c475c69c05281a1d
% canonical-lessons: UR-C499-kapas, UR-C499-dhunna, UR-C499-charkha, UR-C499-katna-spin, UR-C499-ganth

\chapter{Cotton and Thread}
\label{ch:ur-cotton-and-thread}
\hlchaptermodality{499}

\begin{chapteropening}
\textbf{By the end of this chapter:} \emph{I can say how cotton becomes thread.}

Complete the last lesson of chapter 499: I can say how cotton becomes thread.
\end{chapteropening}

\section[\texorpdfstring{\ur{کپاس} (kapās)}{کپاس (kapās)}]{\ur{کپاس} (kapās) — cotton (the plant)}

\label{lesson:UR-C499-kapas}



\begin{quote}
Before the new one: say the Urdu for a large clay pot, then the Urdu for to heat in a fire.
\end{quote}

\subsection*{You'll want to know: \ur{کپاس}}
\textbf{\ur{کپاس}} — \emph{kapās} — \textquotedblleft{}cotton (the plant)\textquotedblright{}.

\begin{quote}
\ur{کسان} \ur{نے} \ur{اس} \ur{سال} \ur{کپاس} \ur{بوئی۔}
\end{quote}

\begin{quote}
\emph{kisān ne is sāl kapās boī.}
\end{quote}

\begin{quote}
The farmer sowed cotton this year.
\end{quote}

Feminine.

\begin{cousinweb}
\textbf{\ur{کپاس}} (\emph{kapās}) is inherited from Sanskrit \emph{karpāsa}, \textquotedblleft{}the cotton plant\textquotedblright{} (CDIAL 2877).
\end{cousinweb}

\begin{grammarlens}[title={What you've built}]
One more word for describing what you lived and felt.
\end{grammarlens}

\subsection*{Your turn}
\begin{itemize}
\raggedright
  \item \emph{Say it:} \emph{kapās}
  \item \emph{Say it:} \emph{kapās}, once more
  \item \emph{Say it:} \emph{kisān ne is sāl kapās boī.} — The farmer sowed cotton this year.
  \item \emph{Recall:} say the Urdu for a large clay pot, then the Urdu for to heat in a fire, then say \emph{kapās} again
\end{itemize}

\begin{tcolorbox}[breakable,colback=teal!4,colframe=teal!35!black,title={Before you move on}]
What does \ur{کپاس} mean? (\textquotedblleft{}Cotton (the plant)\textquotedblright{}.) Say it once more.
\end{tcolorbox}
\section[\texorpdfstring{\ur{دھننا} (dhunnā)}{دھننا (dhunnā)}]{\ur{دھننا} (dhunnā) — to card (cotton)}

\label{lesson:UR-C499-dhunna}



\begin{quote}
Before the new one: say the Urdu for to heat in a fire, then the Urdu for cotton (the plant).
\end{quote}

\subsection*{You'll want to know: \ur{دھننا}}
\textbf{\ur{دھننا}} — \emph{dhunnā} — \textquotedblleft{}to card (cotton)\textquotedblright{}.

\begin{quote}
\ur{دادی} \ur{نے} \ur{رضائی} \ur{کے} \ur{لیے} \ur{کپاس} \ur{دھنی۔}
\end{quote}

\begin{quote}
\emph{dādī ne razāī ke liye kapās dhunī.}
\end{quote}

\begin{quote}
Grandmother carded cotton for the quilt.
\end{quote}

Past, with \ur{نے} (he or she …): \ur{اس} \ur{نے} \ur{کپاس} \ur{دھنی}, \emph{us ne kapās dhunī}.

\begin{cousinweb}
\textbf{\ur{دھننا}} (\emph{dhunnā}) is inherited: CDIAL lists the Hindi cousin \emph{dhunnā}, \textquotedblleft{}to beat, to scutch cotton\textquotedblright{}, under Sanskrit \emph{dhūnōti}, \textquotedblleft{}shakes\textquotedblright{} (CDIAL 6846).
\end{cousinweb}

\begin{grammarlens}[title={What you've built}]
One more word for describing what you lived and felt.
\end{grammarlens}

\subsection*{Your turn}
\begin{itemize}
\raggedright
  \item \emph{Say it:} \emph{dhunnā}
  \item \emph{Say it:} \emph{dhunnā}, once more
  \item \emph{Say it:} \emph{dādī ne razāī ke liye kapās dhunī.} — Grandmother carded cotton for the quilt.
  \item \emph{Recall:} say the Urdu for to heat in a fire, then the Urdu for cotton (the plant), then say \emph{dhunnā} again
\end{itemize}

\begin{tcolorbox}[breakable,colback=teal!4,colframe=teal!35!black,title={Before you move on}]
What does \ur{دھننا} mean? (\textquotedblleft{}To card (cotton)\textquotedblright{}.) Say it once more.
\end{tcolorbox}
\section[\texorpdfstring{\ur{چرخہ} (charkhā)}{چرخہ (charkhā)}]{\ur{چرخہ} (charkhā) — a spinning wheel}

\label{lesson:UR-C499-charkha}



\begin{quote}
Before the new one: say the Urdu for cotton (the plant), then the Urdu for to card (cotton).
\end{quote}

\subsection*{You'll want to know: \ur{چرخہ}}
\textbf{\ur{چرخہ}} — \emph{charkhā} — \textquotedblleft{}a spinning wheel\textquotedblright{}.

\begin{quote}
\ur{نانی} \ur{شام} \ur{کو} \ur{چرخہ} \ur{چلاتی} \ur{تھیں۔}
\end{quote}

\begin{quote}
\emph{nānī shām ko charkhā chalātī thīṅ.}
\end{quote}

\begin{quote}
Grandmother used to work the spinning wheel in the evening.
\end{quote}

Masculine.

\begin{cousinweb}
\textbf{\ur{چرخہ}} (\emph{charkhā}) is usually traced to Persian \emph{charkha}, \textquotedblleft{}a wheel\textquotedblright{}, built on \emph{charkh}, \textquotedblleft{}a wheel; the turning sky\textquotedblright{}.
\end{cousinweb}

\begin{grammarlens}[title={What you've built}]
One more word for describing what you lived and felt.
\end{grammarlens}

\subsection*{Your turn}
\begin{itemize}
\raggedright
  \item \emph{Say it:} \emph{charkhā}
  \item \emph{Say it:} \emph{charkhā}, once more
  \item \emph{Say it:} \emph{nānī shām ko charkhā chalātī thīṅ.} — Grandmother used to work the spinning wheel in the evening.
  \item \emph{Recall:} say the Urdu for cotton (the plant), then the Urdu for to card (cotton), then say \emph{charkhā} again
\end{itemize}

\begin{tcolorbox}[breakable,colback=teal!4,colframe=teal!35!black,title={Before you move on}]
What does \ur{چرخہ} mean? (\textquotedblleft{}A spinning wheel\textquotedblright{}.) Say it once more.
\end{tcolorbox}
\section[\texorpdfstring{\ur{کاتنا} (kātnā)}{کاتنا (kātnā)}]{\ur{کاتنا} (kātnā) — to spin (thread)}

\label{lesson:UR-C499-katna-spin}



\begin{quote}
Before the new one: say the Urdu for to card (cotton), then the Urdu for a spinning wheel.
\end{quote}

\subsection*{You'll want to know: \ur{کاتنا}}
\textbf{\ur{کاتنا}} — \emph{kātnā} — \textquotedblleft{}to spin (thread)\textquotedblright{}.

\begin{quote}
\ur{وہ} \ur{چرخے} \ur{پر} \ur{دھاگا} \ur{کاتتی} \ur{ہے۔}
\end{quote}

\begin{quote}
\emph{voh charkhe par dhāgā kāttī hai.}
\end{quote}

\begin{quote}
She spins thread on the spinning wheel.
\end{quote}

Past, with \ur{نے} (he or she …): \ur{اس} \ur{نے} \ur{دھاگا} \ur{کاتا}, \emph{us ne dhāgā kātā}.

\begin{cousinweb}
\textbf{\ur{کاتنا}} (\emph{kātnā}) is inherited: CDIAL lists the Hindi cousin \emph{kātnā}, \textquotedblleft{}spins\textquotedblright{}, under \emph{*kartati}, a form it reconstructs (CDIAL 2855).
\end{cousinweb}

\begin{grammarlens}[title={What you've built}]
One more word for describing what you lived and felt.
\end{grammarlens}

\subsection*{Your turn}
\begin{itemize}
\raggedright
  \item \emph{Say it:} \emph{kātnā}
  \item \emph{Say it:} \emph{kātnā}, once more
  \item \emph{Say it:} \emph{voh charkhe par dhāgā kāttī hai.} — She spins thread on the spinning wheel.
  \item \emph{Recall:} say the Urdu for to card (cotton), then the Urdu for a spinning wheel, then say \emph{kātnā} again
\end{itemize}

\begin{tcolorbox}[breakable,colback=teal!4,colframe=teal!35!black,title={Before you move on}]
What does \ur{کاتنا} mean? (\textquotedblleft{}To spin (thread)\textquotedblright{}.) Say it once more.
\end{tcolorbox}
\section[\texorpdfstring{\ur{گانٹھ} (gāṅṭh)}{گانٹھ (gāṅṭh)}]{\ur{گانٹھ} (gāṅṭh) — a knot}

\label{lesson:UR-C499-ganth}



\begin{quote}
Before the new one: say the Urdu for a spinning wheel, then the Urdu for to spin (thread).
\end{quote}

\subsection*{You'll want to know: \ur{گانٹھ}}
\textbf{\ur{گانٹھ}} — \emph{gāṅṭh} — \textquotedblleft{}a knot\textquotedblright{}.

\begin{quote}
\ur{اس} \ur{رسی} \ur{میں} \ur{ایک} \ur{گانٹھ} \ur{ہے۔}
\end{quote}

\begin{quote}
\emph{is rassī meṅ ek gāṅṭh hai.}
\end{quote}

\begin{quote}
There is a knot in this rope.
\end{quote}

Feminine.

\begin{cousinweb}
\textbf{\ur{گانٹھ}} (\emph{gāṅṭh}) is inherited from Sanskrit \emph{granthi}, \textquotedblleft{}a knot\textquotedblright{} (CDIAL 4354).
\end{cousinweb}

\begin{grammarlens}[title={What you've built}]
One more word for describing what you lived and felt.
\end{grammarlens}

\subsection*{Your turn}
\begin{itemize}
\raggedright
  \item \emph{Say it:} \emph{gāṅṭh}
  \item \emph{Say it:} \emph{gāṅṭh}, once more
  \item \emph{Say it:} \emph{is rassī meṅ ek gāṅṭh hai.} — There is a knot in this rope.
  \item \emph{Recall:} say the Urdu for a spinning wheel, then the Urdu for to spin (thread), then say \emph{gāṅṭh} again
\end{itemize}

\begin{tcolorbox}[breakable,colback=teal!4,colframe=teal!35!black,title={Before you move on}]
What does \ur{گانٹھ} mean? (\textquotedblleft{}A knot\textquotedblright{}.) Say it once more.
\end{tcolorbox}
